\section{Methodology}\label{sec:method}

\subsection{From the definitions to an implementation}\label{sec:method-overview}

Sections~\ref{sec:SignalSpace} and~\ref{sec:MetricDefinition} define $G$ through three objects that are computed in turn:
the decomposition $\Ev(\varphi,0)$, its canonical form, and the trimmed
canonical decomposition $\Pc(\varphi)$, on which Eq.~\eqref{eq:OWSim} takes a
mean of best matches. Read literally, these definitions are an algorithm, and
we use that algorithm as the \emph{specification} of the metric. It is not how
the metric is computed, because every one of the three objects can be
exponential in the time bound~$T$. Let $\chi = (x \ge 0.2 \wedge x \le 0.4)$.
Then $\Ev(\mathsf{G}_{[0,16]}\mathsf{F}_{[0,4]}\chi, 0)$ is built from $5^{17}$
witness combinations, since every $\mathsf{G}$ instant chooses one of five
$\mathsf{F}$ witnesses; even after identical boxes are merged, about
$1.5\cdot10^{6}$ distinct boxes remain. And
$\canon(\Ev(\mathsf{F}_{[0,20]}\chi,0))$ has $3^{21}-2^{21} \approx 10^{10}$
canonical boxes. Eq.~\eqref{eq:OWSim} would then compare every canonical box of
one formula with every canonical box of the other.

The implementation computes the same number without building any of these
objects. It proceeds in four stages, each run on one formula over the common
grid of Definition~\ref{def:signalspace}:
\[
  \varphi
  \;\xrightarrow{\ \ref{sec:method-diagram}\ }\; \Dg(\varphi)
  \;\xrightarrow{\ \ref{sec:method-canon}\ }\; \Cg(\varphi)
  \;\xrightarrow{\ \ref{sec:method-trim}\ }\; \bigl(\Cg_L(\varphi)\bigr)_{L=0}^{T+1}
  \;\xrightarrow{\ \ref{sec:method-score}\ }\; \OWSim .
\]
The stages are these:
\begin{enumerate}
  \item The signal space $\Sig(\varphi)$ is stored as a hash-consed decision
    diagram $\Dg(\varphi)$, built by the recursion $\Ev$ with a different
    representation.
  \item The runs of Definition~\ref{can:def:run} are read off $\Dg(\varphi)$, and
    relabelling its edges by run gives a diagram $\Cg(\varphi)$ that accepts
    exactly the canonical boxes.
  \item Trimming becomes a partition of $\Cg(\varphi)$ by trimmed length~$L$.
  \item $\OWSim$ is evaluated by a dynamic program that walks both canonical
    diagrams one level at a time.
\end{enumerate}
The final stage is the only one that reads both formulas.
Figure~\ref{fig:pipelines} summarises the four stages. For each one it gives
what the stage replaces in the methodology and the result that proves the two
compute the same object.
Theorem~\ref{thm:compat} states that the result is $\OWSim_{\varphi\to\vartheta}$
of Eq.~\eqref{eq:OWSim}, hence $G$ of Eq.~\eqref{eq:Global}. Everything proved
in Sections~\ref{sec:SignalSpace}--\ref{sec:soundness} and Appendix~\ref{app:A} therefore holds of the computed score, and
Section~\ref{sec:soundness} in particular.

Three properties of the definitional pipeline carry over unchanged:
\begin{itemize}
  \item \textbf{Exact arithmetic.} Every breakpoint is a constant of the formula
    (Lemma~\ref{can:lem:prov}), so atoms, runs and the truncated constraints of
    Eq.~\eqref{eq:PointSimD} all have endpoints in~$\mathbb{Q}$.
  \item \textbf{Negation never complements a region.} Negation is pushed to the
    atoms before evaluation (Section~\ref{can:sec:neg}).
  \item \textbf{No pairwise alignment.} The two formulas share nothing except
    the grid (Section~\ref{sec:align}).
\end{itemize}

\paragraph{Notation.}
Fix the common grid $\Ax = \{0,\dots,T\}\times V$ and list its $n = (T+1)|V|$
axes in time-major order $a_1,\dots,a_n$, with a fixed order on $V$ inside each
instant. We write $t_k$ for the instant of $a_k$. For a variable $v$ let
$\con_v(\varphi)$ be the set of constants $c$ of the atoms $v \bowtie c$ of
$\varphi$.

\subsection{The signal space as a decision diagram}\label{sec:method-diagram}

\paragraph{The fixed arrangement.}
Put $B^*_{(t,v)} = \con_v(\varphi)$ for every instant~$t$. The atoms of axis
$a_k$ over $B^*$ are $\alpha_{k,1},\dots,\alpha_{k,m_k}$ in increasing order,
with $m_k = 2|B^*_{a_k}|+1$, as in Section~\ref{can:sec:arrange}. Unlike $B(\Ev(\varphi,0))$,
this arrangement is known before anything is evaluated, and the diagram needs
exactly that: one fixed alphabet per level.

\begin{definition}[region diagram]\label{def:diagram}
A \emph{node} of level $k\le n$ is a tuple of $m_k$ nodes of level $k+1$, its
\emph{children}; $\nu[i]$ is the child along atom $\alpha_{k,i}$. Level $n+1$
holds the two terminals $\rej$ and $\acc$. The \emph{language} $\Lang(\nu)$ of
a node $\nu$ of level $k$ is the set of atom strings $(i_k,\dots,i_n)$ whose
path from $\nu$ ends in~$\acc$, and a node $\nu$ of level~1 denotes
\[
  \sem{\nu} \;=\; \bigcup_{(i_1,\dots,i_n)\in\Lang(\nu)} \;\prod_{k=1}^{n}\alpha_{k,i_k}
  \;\subseteq\; \mathbb{R}^{\Ax}.
\]
Nodes are only ever created by $\mk(k,\text{children})$, which returns the
existing node when that tuple has been seen before (a unique table). We write
$\rej_k$ and $\acc_k$ for the level-$k$ nodes all of whose children are
$\rej_{k+1}$ and $\acc_{k+1}$, respectively.
\end{definition}

Every edge goes from level $k$ to level $k+1$, so each string of
$\Lang(\nu)$ picks one atom on every axis, and $\sem{\nu}$ is a union of
atom-products. Atom-products partition $\mathbb{R}^{\Ax}$, so distinct strings
denote disjoint, non-empty products, and
\begin{equation}\label{eq:lang-sem}
  \sem{\nu}=\sem{\nu'} \iff \Lang(\nu)=\Lang(\nu').
\end{equation}

\begin{lemma}[canonicity]\label{lem:diag-canonical}
Two nodes of the same level are equal iff their languages are equal. In
particular $\rej_k$ is the only node of level $k$ with empty language.
\end{lemma}

\begin{proof}
($\Rightarrow$) is trivial. ($\Leftarrow$) Induction on $n+1-k$. At level
$n+1$ the terminals have languages $\emptyset$ and $\{()\}$. At level $k\le n$,
$\Lang(\nu[i]) = \{ s : (i)\cdot s \in \Lang(\nu)\}$, so
$\Lang(\nu)=\Lang(\nu')$ gives $\Lang(\nu[i])=\Lang(\nu'[i])$ for every~$i$.
By the hypothesis the children are equal, so the tuples are equal, and the
unique table returns one node for them.
\end{proof}

\paragraph{Boolean operations.}
For $\circ\in\{\wedge,\vee\}$ and nodes $\nu,\nu'$ of level $k$,
$\apply_\circ(\nu,\nu') = \mk\bigl(k, (\apply_\circ(\nu[i],\nu'[i]))_{i}\bigr)$,
with $\apply_\wedge$ and $\apply_\vee$ acting on terminals as conjunction and
disjunction, and results memoised on the pair. By induction on the level,
$\Lang(\apply_\wedge(\nu,\nu')) = \Lang(\nu)\cap\Lang(\nu')$ and
$\Lang(\apply_\vee(\nu,\nu')) = \Lang(\nu)\cup\Lang(\nu')$. So,
by~\eqref{eq:lang-sem} and the partition, $\sem{\apply_\wedge(\nu,\nu')} =
\sem{\nu}\cap\sem{\nu'}$ and $\sem{\apply_\vee(\nu,\nu')} =
\sem{\nu}\cup\sem{\nu'}$. Unlike $\sqcap$ on decompositions, these operations
do not multiply boxes: every node is shared by all the paths that reach it.

\paragraph{Evaluation.}
For $t + h(\varphi) \le T$ define the level-1 node $\Dg(\varphi,t)$ by the
clauses of~$\Ev$ (Section~\ref{can:sec:eval}), with $\sqcap$ and $\cup$ replaced by
$\apply_\wedge$ and $\apply_\vee$:
\begin{align*}
  \Dg(v\bowtie c,t) &= \text{the node accepting exactly the strings whose}\\
  &\phantom{{}={}}\text{$(t,v)$-atom lies in } \{x : x\bowtie c\},\\
  \Dg(\top,t) &= \acc_1,\qquad \Dg(\bot,t)=\rej_1,\qquad
  \Dg(\varphi\wedge\psi,t) = \apply_\wedge\bigl(\Dg(\varphi,t),\Dg(\psi,t)\bigr),\\
  \Dg(\varphi\vee\psi,t) &= \apply_\vee\bigl(\Dg(\varphi,t),\Dg(\psi,t)\bigr),\\
  \Dg(\mathsf{F}_{[a,b]}\varphi,t) &= \textstyle\bigvee_{u=a}^{b}\Dg(\varphi,t+u),\qquad
  \Dg(\mathsf{G}_{[a,b]}\varphi,t) = \textstyle\bigwedge_{u=a}^{b}\Dg(\varphi,t+u),\\
  \Dg(\varphi\,\mathsf{U}_{[a,b]}\,\psi,t) &= \textstyle\bigvee_{u=a}^{b}
     \Bigl(\Dg(\psi,t+u)\wedge\bigwedge_{s=t}^{t+u}\Dg(\varphi,s)\Bigr),
\end{align*}
where $\bigvee$ and $\bigwedge$ fold $\apply_\vee$ from $\rej_1$ and
$\apply_\wedge$ from $\acc_1$. The recursion is memoised on
(subformula, instant). $\Dg(\mathsf{G}_{[0,16]}\mathsf{F}_{[0,4]}\chi,0)$ therefore
evaluates $\Dg(\chi,t)$ once per instant and never forms the $5^{17}$
combinations. Write $\Dg(\varphi) = \Dg(\varphi,0)$.

\begin{lemma}[the diagram denotes the signal space]\label{lem:diag-sem}
Assume \textup{(P1)}. For every $\varphi$ and every $t$ with $t+h(\varphi)\le
T$, $\sem{\Dg(\varphi,t)} = \sem{\Ev(\varphi,t)}$. In particular
$\sem{\Dg(\varphi)} = \Sig(\varphi)$, and $\Sig(\varphi)$ is a union of
$B^*$-atom-products.
\end{lemma}

\begin{proof}
Structural induction, clause by clause against~$\Ev$.

\emph{Atoms.} By Lemma~\ref{can:lem:atom}, $\{x : x\bowtie c\}$ is a union of
intervals with finite endpoints in $\{c\}\subseteq B^*_{(t,v)}$. By
Lemma~\ref{can:lem:dich} each $B^*$-atom of $(t,v)$ therefore lies inside that set
or is disjoint from it. So $\sem{\Dg(v\bowtie c,t)} = \{w : w(t)(v)\bowtie c\}
= \sem{\Ev(v\bowtie c,t)}$.

\emph{Constants.} Immediate, since $\sem{\acc_1}=\mathbb{R}^{\Ax}$ and
$\sem{\rej_1}=\emptyset$.

\emph{Every other clause.} Each is a finite union or intersection of the
arguments of the corresponding $\Ev$ clause, over the same instants. This
includes the closed invariant window $s\in[t,t+u]$ of the until
(Remark~\ref{can:rem:until}). By the paragraph on Boolean operations, $\apply_\vee$
and $\apply_\wedge$ compute exactly that union and intersection, and
Lemma~\ref{can:lem:algebra} says $\cup$ and $\sqcap$ do the same on
decompositions. The two sides therefore agree by the induction hypothesis.

The last claim holds because every $\sem{\nu}$ is a union of atom-products.
\end{proof}

\subsection{Canonicalisation on the diagram}\label{sec:method-canon}

The canonical form (Definition~\ref{can:def:canon}) is defined over the arrangement
$B(P)$ of the decomposition being canonicalised. The implementation works over
$B^*$ instead, which is in general strictly finer: it cuts $x$ at every
constant of $x$ on every instant, including instants where no box of
$\Ev(\varphi,0)$ uses that constant. The first result shows that the canonical
form does not depend on which arrangement is used, provided the region is a
union of its atom-products.

For an arrangement $B'$, write $\fine_{B'}(\Reg)$ for the set of
$B'$-atom-products contained in~$\Reg$. Runs (Definition~\ref{can:def:run}) and
coarse boxes depend on $\Reg$ alone, and $\coarse(X)$ is defined for
$X\in\fine_{B'}(\Reg)$ exactly as in Section~\ref{can:sec:runs}.

\begin{lemma}[arrangement independence]\label{lem:arr-indep}
Let $\Reg\subseteq\mathbb{R}^{\Ax}$ be a union of $B'$-atom-products. Then
\[
  \{\coarse(X) : X\in\fine_{B'}(\Reg)\}
  \;=\; \{C \text{ a coarse box} : \sem{C}\subseteq\Reg\}
  \;=\; \canon(P)
\]
for every decomposition $P$ with $\sem{P}=\Reg$.
\end{lemma}

\begin{proof}
The second equality is Proposition~\ref{can:prop:charac}, applied to $P$. For the
first, we check that the proofs of Lemma~\ref{can:lem:const},
Corollary~\ref{can:cor:runatom}, Lemma~\ref{can:lem:sat} and
Proposition~\ref{can:prop:charac} use the arrangement only through the conclusion
of Lemma~\ref{can:lem:const}. That conclusion also holds for $B'$-atoms:

\emph{Lemma~\ref{can:lem:const} for $B'$.} Let $x,x'$ lie in one $B'$-atom $\alpha$
of axis~$a$, and let $y\in\mathbb{R}^{\Ax\setminus\{a\}}$. The signals $(x,y)$
and $(x',y)$ lie in the same $B'$-atom-product, because they differ only on
$a$, where both lie in~$\alpha$. That product lies inside $\Reg$ or is disjoint
from it, since $\Reg$ is a union of such products. Hence $(x,y)\in\Reg$ iff
$(x',y)\in\Reg$, that is, $\kappa_a(x)=\kappa_a(x')$.

With this in place, Corollary~\ref{can:cor:runatom} gives that every run is a union
of $B'$-atoms. The argument before Definition~\ref{can:def:canon} then makes
$\coarse(X)$ well defined for $X\in\fine_{B'}(\Reg)$. The proofs of
Lemma~\ref{can:lem:sat} and Proposition~\ref{can:prop:charac} go through word for word,
with $\fine_{B'}(\Reg)$ in place of $\fine(P)$, and give the first equality.
\end{proof}

By Lemma~\ref{lem:diag-sem}, $\Sig(\varphi)$ is a union of
$B^*$-atom-products, and $\fine_{B^*}(\Sig(\varphi))$ is exactly the set of
products named by the strings of $\Lang(\Dg(\varphi))$. It remains to compute
the runs from the diagram. Let $R_k$ be the set of nodes of level $k$ reached
from $\Dg(\varphi)$ by a path that avoids the nodes $\rej_j$. By
Lemma~\ref{lem:diag-canonical}, these are the level-$k$ nodes reached by some
prefix of an accepted string.

\begin{lemma}[fibres from the diagram]\label{lem:fibre}
Let $\Sig(\varphi)\neq\emptyset$, let $a=a_k$, and let $\alpha=\alpha_{k,i}$
and $\alpha'=\alpha_{k,i'}$. Then
\begin{enumerate}
  \item $\kappa_a$ is non-empty on $\alpha$ iff $\nu[i]\neq\rej_{k+1}$ for some
    $\nu\in R_k$;
  \item $\kappa_a$ takes the same value on $\alpha$ and $\alpha'$ iff
    $\nu[i]=\nu[i']$ for every $\nu\in R_k$.
\end{enumerate}
\end{lemma}

\begin{proof}
Pick $x\in\alpha$ and split $y\in\mathbb{R}^{\Ax\setminus\{a\}}$ into its part
$y_<$ on $a_1,\dots,a_{k-1}$ and its part $y_>$ on $a_{k+1},\dots,a_n$. There
is exactly one atom string $p$ of length $k-1$ whose product contains $y_<$, and
exactly one string $s$ of length $n-k$ whose product contains $y_>$. Each node
has one child per atom, so $p$ leads from the root to exactly one node $\nu_p$
of level~$k$. Then $(x,y)\in\Sig(\varphi)$ iff $p\cdot(i)\cdot s$ is accepted,
iff $s\in\Lang(\nu_p[i])$. Hence
\[
  \kappa_a(x) \;=\; \bigcup_{p} \;\sem{p}\times\sem{\nu_p[i]},
\]
where $\sem{p}$ and $\sem{\nu_p[i]}$ are the products and unions of products on
the axes before and after~$a$. The sets $\sem{p}$ are non-empty and pairwise
disjoint.

Two such unions are therefore equal iff their terms are equal for every~$p$.
By~\eqref{eq:lang-sem} and Lemma~\ref{lem:diag-canonical}, that holds iff
$\nu_p[i]=\nu_p[i']$ for every~$p$. A prefix with $\nu_p=\rej_k$ has both
children equal to $\rej_{k+1}$, so only $\nu_p\in R_k$ constrains anything,
which gives~(2). Similarly, $\kappa_a(x)\neq\emptyset$ iff some term is
non-empty, iff $\nu_p[i]\neq\rej_{k+1}$ for some $\nu_p\in R_k$, which
gives~(1).
\end{proof}

By Lemma~\ref{lem:arr-indep}, the runs of axis $a_k$ are unions of
$B^*$-atoms. So each run is a maximal block of \emph{consecutive} atoms on
which $\kappa_a$ is non-empty and constant. By Lemma~\ref{lem:fibre}, such a
block is a maximal block of consecutive atoms whose child tuples
$(\nu[i])_{\nu\in R_k}$ are equal and not all~$\rej_{k+1}$. The implementation
scans the atoms in order and extends the current run while this holds. When an
atom with empty $\kappa_a$ is skipped, the next atom no longer touches the
current run: two atoms open at the same breakpoint $b$ are separated by the
missing $[b,b]$. So holes need no special case, exactly as in
Section~\ref{can:sec:runs}.

\begin{proposition}[canonical diagram]\label{prop:canon-diag}
Let $S_{k,1},\dots,S_{k,r_k}$ be the runs of $a_k$, and let $f_k(j)$ be the
index of the first atom of $S_{k,j}$. Define
$\coarse(\rej_k)=\rej'_k$, $\coarse(\acc)=\acc$, and, for $\nu\in R_k$,
\[
  \coarse(\nu) = \mk'\bigl(k,\ (\coarse(\nu[f_k(j)]))_{j=1}^{r_k}\bigr),
\]
where $\mk'$ is the unique table of a second diagram whose level~$k$ has one
edge per run of $a_k$. Let $\Cg(\varphi) = \coarse(\Dg(\varphi))$. Then the
strings of $\Lang(\Cg(\varphi))$ are exactly the canonical boxes of
$\canon(\Ev(\varphi,0))$, one string per canonical box.
\end{proposition}

\begin{proof}
Follow a run string $(j_1,\dots,j_n)$ through $\Cg(\varphi)$. This retraces the
path of the atom string $(f_1(j_1),\dots,f_n(j_n))$ through $\Dg(\varphi)$,
and every node on that path is either some $\rej_k$ or lies in $R_k$. So the
run string is accepted iff $X=\prod_k\alpha_{k,f_k(j_k)}$ lies in
$\fine_{B^*}(\Sig(\varphi))$.

Let $C$ be the coarse box $a_k\mapsto S_{k,j_k}$. If $X\in\fine_{B^*}$, then
$\coarse(X)=C$, and $\sem{C}\subseteq\Sig(\varphi)$ by Lemma~\ref{can:lem:sat}.
Conversely, if $\sem{C}\subseteq\Sig(\varphi)$, then
$X\subseteq\sem{C}\subseteq\Sig(\varphi)$. So the accepted run strings are
exactly the coarse boxes contained in $\Sig(\varphi)$. By
Lemma~\ref{lem:arr-indep} and Lemma~\ref{lem:diag-sem}, these are
$\canon(\Ev(\varphi,0))$. Distinct run strings name distinct coarse boxes.
\end{proof}

Proposition~\ref{prop:canon-diag} uses only Lemma~\ref{can:lem:sat}, not the choice
of $f_k(j)$: by Lemma~\ref{lem:fibre}, any atom of the run has the same child
under every node of $R_k$. $\Cg(\varphi)$ is again built with a unique table,
so Lemma~\ref{lem:diag-canonical} and the operations $\apply_\wedge$ and
$\apply_\vee$ apply to it. The number of canonical boxes, $|\Lang(\Cg(\varphi))|$,
is computed by the memoised recursion $|\Lang(\nu)|=\sum_i|\Lang(\nu[i])|$,
without listing the boxes. For $\mathsf{F}_{[0,20]}\chi$ this yields
$3^{21}-2^{21}$.

\subsection{Trimming as a partition by length}\label{sec:method-trim}

A canonical box $C$ assigns a run $C_a$ to every axis. The constraint is
$\mathtt{undef}$ exactly when $C_a=(-\infty,\infty)$, which is possible only
when $(-\infty,\infty)$ is the sole run of~$a$. Trimming
(Section~\ref{sec:align}) drops the trailing instants at which every
variable is $\mathtt{undef}$, so the trimmed box has time domain
$\{0,\dots,\len(C)-1\}$, where
\[
  \len(C) \;=\; \min\bigl\{L\in\{0,\dots,T+1\} : C_{(t,v)}=(-\infty,\infty)
  \text{ for all } t\ge L,\ v\in V\bigr\}.
\]
For each $L$, let $M_L$ be the diagram over the run alphabet that accepts
exactly the run strings with $\len=L$. It forces the $\mathtt{undef}$ run on
every axis with $t\ge L$ and, if $L\ge1$, requires a run other than
$\mathtt{undef}$ on some axis of instant $L-1$. $M_L$ has at most two nodes per
level. Put $\Cg_L(\varphi)=\apply_\wedge(\Cg(\varphi),M_L)$.

\begin{lemma}[trimming]\label{lem:trim}
Trimming is injective on $\canon(\Ev(\varphi,0))$. The trimmed canonical
decomposition $\Pc(\varphi)$ is therefore the disjoint union over
$L\in\{0,\dots,T+1\}$ of the trimmed boxes of $\Lang(\Cg_L(\varphi))$, and
each box of $\Lang(\Cg_L(\varphi))$ trims to time domain $\{0,\dots,L-1\}$.
\end{lemma}

\begin{proof}
The last two claims follow from the construction of $M_L$ and
Proposition~\ref{prop:canon-diag}. For injectivity, let $C\neq C'$ differ on
axis $(t,v)$. If $t<\min(\len(C),\len(C'))$, both trimmed boxes keep that axis
and differ on it. Otherwise, say $t\ge\len(C)$. Then $C_{(t,v)}=(-\infty,\infty)$,
so $C'_{(t,v)}\neq(-\infty,\infty)$ and $\len(C')>t\ge\len(C)$. The trimmed
boxes then have different time domains.
\end{proof}

In particular $|\Pc(\varphi)|=|\Lang(\Cg(\varphi))|$, the denominator of
Eq.~\eqref{eq:OWSim}.

\subsection{Scoring by dynamic programming}\label{sec:method-score}

Both formulas are canonicalised over the same axes $a_1,\dots,a_n$, so their
diagrams have the same levels, with different run alphabets. Let
$S^{\varphi}_{k,i}$ and $S^{\vartheta}_{k,j}$ be the runs of $a_k$ for the two
formulas. Point similarity depends only on the two constraints compared, so
it is tabulated once per level:
\[
  \pi_k(i,j)=\PointSim\bigl(S^{\varphi}_{k,i},S^{\vartheta}_{k,j}\bigr),
\]
with the $D$-truncation and the endpoint set $E(c_1,c_2)$ of
Definitions~\ref{def:Dimension} and~\ref{def:epsmeasure}. For canonical boxes
$\beta_1=(i_1,\dots,i_n)$ of length $L_1$ and $\beta_2=(j_1,\dots,j_n)$ of
length $L_2$, the trimmed time domains are $\{0,\dots,L_1-1\}$ and
$\{0,\dots,L_2-1\}$, so Eq.~\eqref{eq:BoxSim} reads
\begin{equation}\label{eq:boxsim-len}
\begin{aligned}
  \BoxSim(\beta_1,\beta_2) &= \frac{W_{L_1,L_2}(\beta_1,\beta_2)}{\max(L_1,L_2)\,|V|},\\
  W_{L_1,L_2}(\beta_1,\beta_2) &= \sum_{k\,:\,t_k<\min(L_1,L_2)}\pi_k(i_k,j_k),
\end{aligned}
\end{equation}
and $\BoxSim=1$ when $L_1=L_2=0$. The normalisation depends on $\beta_2$ only
through $L_2$. The best match of $\beta_1$ is therefore the best, over lengths
$L_2$ that occur in $\Pc(\vartheta)$, of the length-$L_2$ normalisation applied
to
\[
  W^*_{L_2}(\beta_1)=\max_{\beta_2\in\Lang(\Cg_{L_2}(\vartheta))}W_{L_1,L_2}(\beta_1,\beta_2).
\]

\paragraph{The program.}
Fix $L_1$. A \emph{state} after $k$ levels is a pair $(\mu, w)$:
\begin{itemize}
  \item $\mu$ is the node of $\Cg_{L_1}(\varphi)$ reached by a prefix
    $(i_1,\dots,i_k)$ of a box of $\varphi$;
  \item $w$ is a \emph{Viterbi vector}. It maps each pair $(L_2,\nu)$, with
    $\nu\neq\rej_{k+1}$ a node of $\Cg_{L_2}(\vartheta)$ at level $k+1$, to the
    largest partial sum $\sum_{k'\le k,\ t_{k'}<\min(L_1,L_2)}\pi_{k'}(i_{k'},j_{k'})$
    over the prefixes $(j_1,\dots,j_k)$ that lead to~$\nu$.
\end{itemize}
The initial state is $\bigl(\Cg_{L_1}(\varphi), \{(L_2,\Cg_{L_2}(\vartheta))\mapsto 0\}_{L_2}\bigr)$.
One step along atom $i$ of $\mu$, with $\mu[i]\neq\rej$, goes to
$(\mu[i],w')$, where
\begin{equation}\label{eq:viterbi}
  w'(L_2,\nu') \;=\; \max_{\substack{(L_2,\nu)\in\mathrm{dom}\,w,\ j\,:\\ \nu[j]=\nu'}}
  \Bigl(w(L_2,\nu) + [\,t_{k+1}<\min(L_1,L_2)\,]\;\pi_{k+1}(i,j)\Bigr).
\end{equation}
States are kept in a table with multiplicities, and two prefixes that reach the
same state are merged by adding their counts. After $n$ levels, every
surviving state has $\mu=\acc$, and its vector is defined on $(L_2,\acc)$ for
each $L_2$ that occurs in $\Pc(\vartheta)$.

\begin{theorem}[compatibility]\label{thm:compat}
Let $D$ be well formed (Definition~\ref{def:Dimension}) and $\varepsilon>0$.
Suppose $\Pc(\varphi)$ and $\Pc(\vartheta)$ are both non-empty. For every $L_1$,
sum over the final states $(\acc,w)$, with multiplicity $c$, the quantity
\[
  c\cdot\max_{L_2}\ \begin{cases}1 & \text{if } L_1=L_2=0,\\[2pt]
  \dfrac{w(L_2,\acc)}{\max(L_1,L_2)\,|V|} & \text{otherwise.}\end{cases}
\]
The sum of these totals over $L_1$, divided by $|\Lang(\Cg(\varphi))|$, equals
$\OWSim_{\varphi\to\vartheta}$ of Eq.~\eqref{eq:OWSim}. The empty cases of
Eq.~\eqref{eq:OWSim} are decided directly from $|\Lang(\Cg(\cdot))|$. Hence
the score computed from the two one-way values by Eq.~\eqref{eq:Global} is
$G(\varphi,\vartheta)$.
\end{theorem}

\begin{proof}
There are three claims.

\emph{(i) The Viterbi vector is exact.} Fix a prefix $(i_1,\dots,i_k)$ of a box
of $\varphi$. We show by induction on $k$ that $w(L_2,\nu)$ is the stated
maximum. For $k=0$ the empty prefix leads to the root with sum~$0$. For the
step, each node has one child per atom, so a prefix of length $k+1$ leading to
$\nu'$ is uniquely a prefix of length $k$ leading to some $\nu$, followed by an
atom $j$ with $\nu[j]=\nu'$. Its partial sum is the sum of the shorter prefix
plus the $(k+1)$-st term. The maximum over all such prefixes is therefore the
maximum over $(\nu,j)$ of [the best sum at $\nu$] plus that term, which
is~\eqref{eq:viterbi}. Prefixes leading to $\rej$ have no accepted completion
(Lemma~\ref{lem:diag-canonical}) and can be dropped. At level $n$ the only
node left is $\acc$, so $w(L_2,\acc)=W^*_{L_2}(\beta_1)$.

\emph{(ii) Merging states is exact.} The successor of a state along atom $i$ at
level $k+1$ is determined by $(\mu,w,i)$ alone, since the indicator
in~\eqref{eq:viterbi} depends only on $k$, $L_1$ and $L_2$. Two prefixes that
reach the same state therefore have the same successors, and in the end the
same final vector. Adding their counts replaces a sum of equal terms by a
product. So the total for $L_1$ equals
\[
  \sum_{\beta_1\in\Lang(\Cg_{L_1}(\varphi))}\;\max_{\beta_2\in\Pc(\vartheta)}\BoxSim(\beta_1,\beta_2)
\]
by~(i) and~\eqref{eq:boxsim-len}.

\emph{(iii) Summing over lengths gives $\OWSim$.} As $L_1$ ranges over
$\{0,\dots,T+1\}$, the languages $\Lang(\Cg_{L_1}(\varphi))$ partition
$\Pc(\varphi)$ by Lemma~\ref{lem:trim}, so the sum over $L_1$ runs over all of $\Pc(\varphi)$.
Also $|\Pc(\varphi)|=|\Lang(\Cg(\varphi))|$. This is the third case of
Eq.~\eqref{eq:OWSim}. The first two cases hold because
$\Pc(\cdot)=\emptyset$ iff $|\Lang(\Cg(\cdot))|=0$.
\end{proof}

\begin{corollary}\label{cor:compat}
For every well-formed $D$ and every $\varepsilon>0$, the computed score equals
$G(\varphi,\vartheta)$ of Eq.~\eqref{eq:Global}. In particular it equals~1 iff
$\varphi\equiv\vartheta$ (Section~\ref{sec:soundness}).
\end{corollary}

\begin{remark}[what is exact, and what the implementation checks]\label{rem:impl}
Breakpoints, runs and truncated constraints are exact rationals (for rational
$D$ and $\varepsilon$), and each ratio $\pi_k(i,j)$ is computed exactly and
then rounded once to floating point. The sums in~\eqref{eq:viterbi} and in
Theorem~\ref{thm:compat} are accumulated in floating point, in a different
order from a box-by-box evaluation of Eq.~\eqref{eq:OWSim}. The two therefore
agree up to rounding, and Theorem~\ref{thm:compat} is a statement about exact
arithmetic.

The preconditions of the metric are enforced rather than assumed. A $D$ that is
not well formed is refused, not clamped, because clamping would report a score
for a $D$ the user did not ask for. $\varepsilon\le0$ is refused because
Lemma~\ref{lem:pointsim-identity} needs $\varepsilon>0$: with $\varepsilon=0$,
$x>0$ and $x\ge0$ would score~1.

As an independent check, the implementation is compared against the literal
pipeline $\Ev\to\canon\to$ trim $\to$ Eq.~\eqref{eq:OWSim} wherever the latter
terminates. The comparison covers the number of canonical boxes, the score on
fixed pairs, and the score on randomly generated pairs of formulas and their
equivalent rewrites.
\end{remark}

\begin{remark}[cost]\label{rem:cost}
Evaluation costs one $\apply$ per clause instance, each bounded by the product
of the sizes of its arguments. Canonicalisation is linear in the size of
$\Dg(\varphi)$. Scoring costs, per level, the number of distinct states times
the size of their vectors. None of these quantities is bounded by the number of
boxes or canonical boxes. On the families of Section~\ref{sec:results}
they remain small: for instance, the pair
$\mathsf{F}_{[0,20]}\chi$ and $\mathsf{G}_{[0,16]}\mathsf{F}_{[0,4]}\chi$ is scored
without enumerating its $\approx10^{10}$ canonical boxes. We claim no
polynomial bound, however: the number of distinct Viterbi vectors can grow
with the horizon.
\end{remark}

\begin{remark}[reuse across comparisons]
$\Cg(\varphi)$ depends only on $\varphi$ and the grid $(T,V)$, not on the other
formula. When many formulas are compared pairwise over one grid, each diagram
is built once, and only the dynamic program of this section is run per pair.
\end{remark}
